\section*{参考文献}\srcpage{53}

\begin{list}{}{\setlength{\leftmargin}{2em}\setlength{\itemindent}{-2em}\setlength{\labelwidth}{0pt}\setlength{\labelsep}{0pt}\setlength{\itemsep}{0.3em}\setlength{\parsep}{0pt}\setlength{\topsep}{0.4em}}

\item Avraham Aizenbud and Nir Avni (2016). “Representation growth and rational singularities of the moduli space of local systems”, \textit{Invent. Math.} \textbf{204} (1), pp. 245–316.
\item Dima Arinkin, Dario Beraldo, Justin Campbell, Lin Chen, Joakim Færgeman, Dennis Gaitsgory, Kevin Lin, Sam Raskin, and Nick Rozenblyum (2024). \textit{Proof of the geometric Langlands conjecture II: Kac-Moody localization and the FLE}. arXiv: 2405.03648.
\item Dima Arinkin, Dario Beraldo, Lin Chen, Joakim Færgeman, Dennis Gaitsgory, Kevin Lin, Sam Raskin, and Nick Rozenblyum (2024). \textit{Proof of the geometric Langlands conjecture IV: ambidexterity}. arXiv: 2409.08670.
\item Dima Arinkin and Dennis Gaitsgory (2015). “Singular support of coherent sheaves and the geometric Langlands conjecture”, \textit{Selecta Math. (N.S.)} \textbf{21} (1), pp. 1–199.
\item Dima Arinkin, Dennis Gaitsgory, David Kazhdan, Sam Raskin, Nick Rozenblyum, and Yakov Varshavsky (2020a). \textit{Duality for automorphic sheaves with nilpotent singular support}. arXiv: 2012.07665.
\srcpage{54}
\item Dima Arinkin, Dennis Gaitsgory, David Kazhdan, Sam Raskin, Nick Rozenblyum, and Yakov Varshavsky (2020b). \textit{The stack of local systems with restricted variation and geometric Langlands theory with nilpotent singular support}. arXiv: 2010.01906.
\item ——— (2021). \textit{Automorphic functions as the trace of Frobenius}. arXiv: 2102.07906.
\item James Arthur (1989). “Unipotent automorphic representations: conjectures”, in: \textit{Orbites unipotentes et représentations, II}. Astérisque 171-172, pp. 13–71.
\item Alexander Beilinson (2016). “Constructible sheaves are holonomic”, \textit{Selecta Math. (N.S.)} \textbf{22} (4), pp. 1797–1819.
\item Alexander Beilinson and Vladimir Drinfeld (1996). “Quantization of Hitchin's fibration and Langlands' program”, in: \textit{Algebraic and geometric methods in mathematical physics (Kaciveli, 1993)}. Math. Phys. Stud. (Kluwer Acad. Publ.), vol. 19, pp. 3–7.
\item David Ben-Zvi (2026). “What is the geometric Langlands correspondence about?”, \textit{Current Events Bulletin of the AMS}, pp. 4–36.
\item David Ben-Zvi and David Nadler (2013). “Loop spaces and representations”, \textit{Duke Math. J.} \textbf{162} (9), pp. 1587–1619.
\item ——— (2018). “Betti geometric Langlands”, in: \textit{Algebraic geometry: Salt Lake City 2015}. Proc. Sympos. Pure Math. (AMS), vol. 97.2, pp. 3–41.
\item Dario Beraldo, David Kazhdan, and Tomer M. Schlank (2022). \textit{Contractibility of the space of generic opers for classical groups}. arXiv: 1801.00655v2.
\item Kevin Buzzard and Toby Gee (2014). “The conjectural connections between automorphic representations and Galois representations”, in: \textit{Automorphic forms and Galois representations. Vol. 1}. London Math. Soc. Lecture Note Ser. (Cambridge Univ. Press), vol. 414, pp. 135–187.
\item Justin Campbell, Lin Chen, Joakim Færgeman, Dennis Gaitsgory, Kevin Lin, Sam Raskin, and Nick Rozenblyum (2024). \textit{Proof of the geometric Langlands conjecture III: compatibility with parabolic induction}. arXiv: 2409.07051.
\item Vladimir Drinfeld (1983). “Two-dimensional $l$-adic representations of the fundamental group of a curve over a finite field and automorphic forms on $GL(2)$”, \textit{Amer. J. Math.} \textbf{105} (1), pp. 85–114.
\item Vladimir Drinfeld and Dennis Gaitsgory (2016). “Geometric constant term functor(s)”, \textit{Selecta Math. (N.S.)} \textbf{22} (4), pp. 1881–1951.
\item Alexander Efimov (2025). \textit{K-theory and localizing invariants of large categories}. arXiv: 2405.12169v3.
\item Matthew Emerton, Toby Gee, and Eugen Hellmann (2025). \textit{An introduction to the categorical p-adic Langlands program}. arXiv: 2210.01404v5.
\item Joakim Færgeman and Sam Raskin (2025). “Non-vanishing of geometric Whittaker coefficients for reductive groups”, \textit{J. Amer. Math. Soc.} \textbf{38} (4), pp. 919–995.
\item Laurent Fargues (2025). “Geometrization of the local Langlands correspondence: an overview”, \textit{J. Math. Sci. Univ. Tokyo} \textbf{32} (2), pp. 157–240.
\item Laurent Fargues and Peter Scholze (2024). \textit{Geometrization of the local Langlands correspondence}, to appear in Astérisque. arXiv: 2102.13459.
\srcpage{55}
\item Boris Feigin and Edward Frenkel (1992). “Affine Kac-Moody algebras at the critical level and Gelfand-Dikiĭ algebras”, in: \textit{Infinite analysis, Part A, B (Kyoto, 1991)}. Adv. Ser. Math. Phys. (World Sci. Publ.), vol. 16, pp. 197–215.
\item Javier Fresán and Peter Jossen (2025). “Exponential motives”. \url{https://javier.fresan.perso.math.cnrs.fr/expmot.pdf}.
\item Dennis Gaitsgory (2015). “Outline of the proof of the geometric Langlands conjecture for $GL_2$”, \textit{Astérisque} (370), pp. 1–112.
\item Dennis Gaitsgory and Sam Raskin (2025a). \textit{Geometric Langlands in positive characteristic from characteristic zero}. arXiv: 2508.02237.
\item ——— (2025b). \textit{Proof of the geometric Langlands conjecture I: construction of the functor}. arXiv: 2405.03599v3.
\item ——— (2026). \textit{Proof of the geometric Langlands conjecture V: the multiplicity one theorem}. arXiv: 2409.09856v3.
\item Dennis Gaitsgory and Nick Rozenblyum (2017). \textit{A study in derived algebraic geometry. Vol. I. Correspondences and duality}. Math. Surveys and Monographs (AMS), vol. 221.
\item Martin Gallauer and Simon Pepin Lehalleur (2022). “Exponentiation of coefficient systems and exponential motives”. arXiv: 2211.17247.
\item Linus Hamann, David Hansen, and Peter Scholze (2024). \textit{Geometric Eisenstein series I: finiteness theorems}. arXiv: 2409.07363.
\item David Hansen and Peter Scholze (2023). “Relative perversity”, \textit{Commun. Am. Math. Soc.} \textbf{3}, pp. 631–668.
\item Claudius Heyer and Lucas Mann (2024). \textit{6-Functor Formalisms and Smooth Representations}. arXiv: 2410.13038.
\item Tim Kuppel (2024). \textit{Proper Sheaves for Schemes and Stacks}. Master Thesis, University of Bonn.
\item Vincent Lafforgue (2018). “Chtoucas pour les groupes réductifs et paramétrisation de Langlands globale”, \textit{J. Amer. Math. Soc.} \textbf{31} (3), pp. 719–891.
\item Vincent Lafforgue and Xinwen Zhu (2019). \textit{Décomposition au-dessus des paramètres de Langlands elliptiques}. arXiv: 1811.07976v2.
\item Gérard Laumon (1987). “Correspondance de Langlands géométrique pour les corps de fonctions”, \textit{Duke Math. J.} \textbf{54} (2), pp. 309–359.
\item ——— (1996). \textit{Transformation de Fourier géneralisée}. arXiv: alg-geom/9603004.
\item Yifeng Liu and Weizhe Zheng (2024). \textit{Enhanced six operations and base change theorem for higher Artin stacks}. arXiv: 1211.5948.
\item Jacob Lurie (2017). “Higher Algebra”. \url{https://www.math.ias.edu/~lurie/papers/HA.pdf}.
\item Lucas Mann (2022). \textit{A p-Adic 6-Functor Formalism in Rigid-Analytic Geometry}. arXiv: 2206.02022.
\item Ivan Mirković and Kari Vilonen (2007). “Geometric Langlands duality and representations of algebraic groups over commutative rings”, \textit{Ann. of Math. (2)} \textbf{166} (1), pp. 95–143.
\srcpage{56}
\item David Nadler and Jeremy Taylor (2025). “The Whittaker functional is a shifted microstalk”, \textit{Transform. Groups} \textbf{30} (3), pp. 1425–1450.
\item David Nadler and Zhiwei Yun (2019). “Spectral action in Betti geometric Langlands”, \textit{Israel J. Math.} \textbf{232} (1), pp. 299–349.
\item Sam Raskin (2025). “Conjectures of Arthur and Ramanujan for unramified automorphic forms: announcement”. \url{https://www.samraskin.net/icm.pdf}.
\item Timo Richarz and Jakob Scholbach (2021). “The motivic Satake equivalence”, \textit{Math. Ann.} \textbf{380} (3-4), pp. 1595–1653.
\item Peter Scholze (2024). \textit{Geometrization of the real local Langlands correspondence}.
\item ——— (2025). “Wild Betti sheaves”, in: \textit{Algebraic Geometry and the Langlands Program, a volume in honour of Gérard Laumon}.
\item ——— (2026). “Geometry and Higher Category Theory”. \url{https://people.mpim-bonn.mpg.de/scholze/Gestalten.pdf}.
\item Dmitry Tamarkin (2018). “Microlocal condition for non-displaceability”, in: \textit{Algebraic and analytic microlocal analysis}. Springer Proc. Math. Stat. vol. 269, pp. 99–223.
\item Thibaud van den Hove (2024). \textit{The integral motivic Satake equivalence for ramified groups}. arXiv: 2404.15694.
\item André Weil (1974). \textit{Basic number theory}. Third edition. Die Grundlehren der mathematischen Wissenschaften (Springer), Bd. 144.
\item Cong Xue (2025). “Smoothness of cohomology sheaves of stacks of shtukas”, in: \textit{Algebraic Geometry and the Langlands Program, a volume in honour of Gérard Laumon}.
\item Xinwen Zhu (2024). “A note on integral Satake isomorphisms”, in: \textit{Arithmetic geometry}. Tata Inst. Fundam. Res. Stud. Math, vol. 41, pp. 469–489.
\item ——— (2025). “Coherent sheaves on the stack of Langlands parameters”, in: \textit{The Langlands program. II—geometrization of the Langlands correspondence}. Proc. Sympos. Pure Math. (AMS), vol. 112, pp. 39–123.


\end{list}

\bigskip
\begin{flushright}
Peter Scholze\\
Max-Planck-Institut für Mathematik\\
Vivatsgasse 7\\
53111 Bonn\\
\textit{E-mail} : scholze@mpim-bonn.mpg.de
\end{flushright}
